\documentclass[11pt]{article}

% ── Packages ────────────────────────────────────────────────────────────────
\usepackage[margin=1in]{geometry}
\usepackage{amsmath, amssymb, amsthm}
\usepackage{mathtools}
\usepackage{enumitem}
\usepackage{hyperref}
\usepackage{parskip}

% ── Theorem environments ─────────────────────────────────────────────────────
\newtheorem{theorem}{Theorem}[section]
\newtheorem{lemma}[theorem]{Lemma}
\newtheorem{corollary}[theorem]{Corollary}
\newtheorem{proposition}[theorem]{Proposition}
\theoremstyle{definition}
\newtheorem{definition}[theorem]{Definition}
\newtheorem{example}[theorem]{Example}
\theoremstyle{remark}
\newtheorem{remark}[theorem]{Remark}

% ── Custom commands ───────────────────────────────────────────────────────────
\newcommand{\R}{\mathbb{R}}
\newcommand{\E}{\mathbb{E}}
\newcommand{\norm}[1]{\left\lVert#1\right\rVert}
\newcommand{\inner}[2]{\langle #1,\, #2 \rangle}
\newcommand{\T}{^{\top}}

% ── Document ─────────────────────────────────────────────────────────────────
\begin{document}

\title{\textbf{Regression: Learning}\\[0.4em]
       \large Notes Template}
\author{}
\date{}
\maketitle
\tableofcontents
\newpage

% ─────────────────────────────────────────────────────────────────────────────
\section{Problem}
% ─────────────────────────────────────────────────────────────────────────────

\subsection{Setup}

% TODO: Describe the regression learning problem here.
% Example prompts:
%   - What is the hypothesis class?
%   - What training data / likelihood model is assumed?
%   - What optimisation criterion is used (MLE, MAP, ERM, …)?

\begin{definition}[Regression Learning Problem]
    % TODO: Fill in the formal definition.
    Given a training dataset
    $\mathcal{D} = \{(x^{(i)}, y^{(i)})\}_{i=1}^{n}$
    with $x^{(i)} \in \R^{d}$ and $y^{(i)} \in \R$,
    the learning problem is to find a parameter vector $\theta$ such that \ldots
\end{definition}

\subsection{Assumptions}

\begin{enumerate}[label=\textbf{A\arabic*.}]
    \item % TODO: State assumption 1 (e.g., noise model).
    \item % TODO: State assumption 2 (e.g., model family).
    \item % TODO: Add further assumptions as needed.
\end{enumerate}

\subsection{Objective}

% TODO: State the learning objective (e.g., empirical risk, negative log-likelihood).
\[
    % TODO: Replace with the actual objective.
    \hat{\theta} \;=\; \arg\min_{\theta} \; \frac{1}{n}\sum_{i=1}^{n}
        \mathcal{L}\!\left(y^{(i)},\; f_{\theta}(x^{(i)})\right)
\]

% ─────────────────────────────────────────────────────────────────────────────
\section{Derivation}
% ─────────────────────────────────────────────────────────────────────────────

The following derivation establishes the optimal learning algorithm step by step.

% ── Step 1 ────────────────────────────────────────────────────────────────────
\subsection{Step 1: [Title]}

% TODO: Describe what this step accomplishes (e.g., write out the likelihood).

\begin{proof}
    % TODO: Fill in the proof for Step 1.
    \[
        \text{(derivation goes here)}
    \]
\end{proof}

% ── Step 2 ────────────────────────────────────────────────────────────────────
\subsection{Step 2: [Title]}

% TODO: Describe what this step accomplishes (e.g., take the log, expand squares).

\begin{proof}
    % TODO: Fill in the proof for Step 2.
    \begin{align}
        &\text{(left-hand side)} \notag \\
        &= \text{(right-hand side)} \label{eq:step2}
    \end{align}
\end{proof}

% ── Step 3 ────────────────────────────────────────────────────────────────────
\subsection{Step 3: [Title]}

% TODO: Describe what this step accomplishes (e.g., take gradient and set to zero).

\begin{proof}
    % TODO: Fill in the proof for Step 3.
    \[
        \nabla_{\theta}\; \mathcal{L}(\theta) \;=\; \mathbf{0}
        \quad \Longrightarrow \quad
        \text{(closed-form solution or gradient update)}
    \]
\end{proof}

% ── Step 4 ────────────────────────────────────────────────────────────────────
\subsection{Step 4: [Title]}

% TODO: Describe what this step accomplishes (e.g., verify the solution is a minimum).

\begin{proof}
    % TODO: Fill in the proof for Step 4.
    \[
        \text{(derivation goes here)}
    \]
\end{proof}

% ── Final Result ──────────────────────────────────────────────────────────────
\subsection{Final Result}

\begin{theorem}[Optimal Regression Learner]
    % TODO: State the final theorem.
    Under assumptions \textbf{A1}--\textbf{A3}, the optimal parameter estimate is
    \[
        % TODO: Replace with the actual result (e.g., OLS, ridge solution).
        \hat{\theta}^{*} \;=\; \left(X\T X\right)^{-1} X\T y.
    \]
\end{theorem}

\begin{proof}
    Follows directly from Steps 1--4. \qed
\end{proof}

% ─────────────────────────────────────────────────────────────────────────────
\end{document}
